% ------------------------------------------------------------------------------
% Chapter 4
% ------------------------------------------------------------------------------
\chapter{Results} 
\label{cha:results}

Overall, the DUT was verified to be functionally correct.
Both the formal verification tool and the UVM architecture confirm this result.
However, the code coverage analysis revealed some minor structural inefficiencies in the implementation.
While these do not endanger the functionality of the DUT, they might lead to unoptimized results upon synthesis.

\section{Coverage Analysis}

From the coverage report, we can observe that the deterministic test vectors completely cover the functional specification (100\%).
To clarify the functional metrics:
\begin{itemize}
    \item \texttt{COVERGROUP} refers to the UVM coverage blocks.
    \item \texttt{DIRECTIVE} refers to the two cover properties in the Hardware coverage.
    \item \texttt{ASSERTION} refers to the Hardware-based assertions.
\end{itemize}

The Total Coverage score calculates the average of these three functional coverage categories along with the seven code coverage points (detailed below in Table \ref{tab:code_coverage}).

\begin{verbatim}
TOTAL COVERGROUP COVERAGE: 100.00%  COVERGROUP TYPES: 2
TOTAL DIRECTIVE COVERAGE: 100.00%  COVERS: 2
TOTAL ASSERTION COVERAGE: 100.00%  ASSERTIONS: 7
Total Coverage By Instance (filtered view): 94.68%
\end{verbatim}

The Total Coverage score is slightly reduced by a few specific structural loopholes identified during the code coverage extraction:

\begin{itemize}
    \item \textbf{Branch [95.45\%]}: Contains 1 uncovered bin.
    \begin{itemize}
        \item \textit{Line 53}: A branch is never evaluated as false.
        This is the consequence of a concatenation of conditions that makes the false outcome logically unreachable.
        The \texttt{if} statement could theoretically be removed.
    \end{itemize}
    
    \item \textbf{Condition [71.42\%]}: Contains 2 uncovered bins.
    \begin{itemize}
        \item \textit{Line 53}: The condition is never false.
        Similar to the branch issue, previous conditions make the false outcome unreachable, making the statement redundant.
        \item \textit{Line 75}: Part of the condition (\texttt{in\_ready}) is never false.
        This happens because a preceding condition (\texttt{state == IDLE}) already forces this signal high, making the false evaluation unreachable.
    \end{itemize}
    
    \item \textbf{FSM Transitions [80.00\%]}: Contains 1 uncovered bin.
    \begin{itemize}
        \item \textit{Line 63}: The state machine never transitions directly from \texttt{RUN} to \texttt{IDLE}.
        This occurs because the reset signal was only exercised during the initial initialization of the DUT, and an active reset is the only mechanism capable of producing this specific mid-transaction transition.
    \end{itemize}
\end{itemize}

\begin{table}[H]
    \centering
    \begin{tabular}{|l|c|c|c|c|}
        \toprule
        \textbf{Enabled Coverage} & \textbf{Bins} & \textbf{Hits} & \textbf{Misses} & \textbf{Coverage} \\ \midrule
        Branches        & 22  & 21  & 1 & 95.45\%  \\ 
        Conditions      & 7   & 5   & 2 & 71.42\%  \\ 
        Expressions     & 2   & 2   & 0 & 100.00\% \\ 
        FSM States      & 3   & 3   & 0 & 100.00\% \\ 
        FSM Transitions & 5   & 4   & 1 & 80.00\%  \\ 
        Statements      & 25  & 25  & 0 & 100.00\% \\ 
        Toggles         & 527 & 527 & 0 & 100.00\% \\ \bottomrule
    \end{tabular}
    \caption{Structural Code Coverage Summary}
    \label{tab:code_coverage}
\end{table}

In the end we consider the code coverage complete as FSM transitions are completely covered by the Formal verification and the other uncovered bins are unreachable code.

\section{Bug}

I introduced a bug in the DUT that should give wrong results as the formula is completely wrong.
We expect the DUT to give wrong gcd for operands that do not fall in the early exit conditions.
Here is the code diff: 
\begin{lstlisting}[language=Verilog, escapeinside={(*}{*)}]
always_comb begin
    a_next = a_reg;
    b_next = b_reg;
    if (state == RUN) begin
(*\textcolor{red}{- \quad\quad\quad\quad if (a\_reg \textgreater b\_reg)}*)
(*\textcolor{green!60!black}{+ \quad\quad\quad\quad if (a\_reg \textless b\_reg)}*)
            a\_next = a\_reg - b\_reg
(*\textcolor{red}{- \quad\quad\quad\quad else if (b\_reg \textgreater a\_reg)}*)
(*\textcolor{green!60!black}{+ \quad\quad\quad\quad else if (b\_reg \textless a\_reg)}*)
           b\_next = b\_reg - a\_reg;
    end
end
\end{lstlisting}

\subsection*{Analysis}


